\section{Introducción}

Los Sistemas de Transporte Inteligentes (ITS, por sus siglas en inglés) y el monitoreo del comportamiento de conductores representan pilares fundamentales en la evolución de las ciudades inteligentes y la seguridad vial global. Ante el aumento en el riesgo de siniestros, resulta crucial desarrollar soluciones tecnológicas automatizadas, no invasivas y en tiempo real para supervisar y regular las conductas viales \cite{kumar2026intelligent, ahmad2026comparative}. Tradicionalmente dependientes de auditorías manuales o de sensores físicos de hardware invasivos y de alto costo (como bucles de inducción o radares), el advenimiento del procesamiento digital de imágenes y el aprendizaje profundo ha propiciado un cambio de paradigma hacia sistemas inteligentes basados en cámaras ópticas que procesan video en tiempo real de forma inteligente \cite{kumar2026uavbased, chen2026hybrid}.

En la literatura científica reciente, esta transición tecnológica se ha bifurcado principalmente en dos grandes vertientes. Por un lado, el monitoreo de conductores se enfoca en mitigar factores conductuales directos (fatiga, somnolencia y distracción al volante) mediante cámaras instaladas en el tablero (\textit{dashcams}). Investigaciones recientes como \cite{liu2026research, abarghooei2026driving, hu2026trustworthy, srivastava2026traffic, pavan2026aipowered} utilizan redes neuronales convolucionales y modelos Transformer para capturar características espaciotemporales de la fatiga, estimar variables fisiológicas (apertura ocular PERCLOS, parpadeo, rotación de la cabeza) y emitir alertas acústicas o vibratorias antes de una colisión inminente. Por otro lado, la regulación activa del tráfico urbano y la penalización automatizada de infracciones (como cruce de semáforos en rojo, invasión de carriles preferenciales o exceso de velocidad) ha migrado a sistemas inteligentes basados en cámaras de CCTV y drones \cite{zia2025advancing, bhavsar2025evaluating}. Estudios como \cite{zambrano2026monitoring, rahman2026costeffective, l2026safestreets} confirman la efectividad de detectores de objetos de una sola etapa (familia YOLO) y algoritmos de seguimiento de objetos múltiples (MOT) para automatizar este control de manera continua y recopilar flujos vehiculares agregados útiles en la gestión del tráfico.

Sin embargo, el despliegue de estas tecnologías en escenarios reales del mundo físico enfrenta desafíos técnicos y operativos severos que constituyen el planteamiento del problema. Primero, la variabilidad lumínica extrema (reflejos, conducción nocturna) y las oclusiones (anteojos o mascarillas en conductores; vehículos grandes bloqueando la visibilidad vial) degradan la precisión de los modelos. Segundo, las arquitecturas de deep learning más robustas conllevan una alta demanda computacional y energética, lo que dificulta su ejecución a alta velocidad y bajo consumo en hardware embebido en el borde (\textit{Edge Computing}). Tercero, el sesgo de los conjuntos de datos, recopilados mayormente en entornos controlados, limita la generalización hacia condiciones de conducción naturalista real. Finalmente, la captura biométrica continua genera problemas de privacidad que exigen procesamiento local o aprendizaje federado. Esto define la brecha técnica entre el desempeño ideal de los modelos y las restricciones operativas reales en carretera.

Bajo este contexto, esta revisión sistemática se justifica por la necesidad de consolidar y evaluar cuantitativamente la copiosa y dispersa literatura científica publicada en el periodo 2020--2026. Ante la irrupción exponencial de nuevos algoritmos y plataformas de hardware, este trabajo proporciona un marco de referencia que estructura el campo bajo los estándares PRISMA y PICO para identificar configuraciones metodológicas y tecnológicas idóneas según el entorno, analizando métricas paramétricas de rendimiento (precisión, mAP, FPS) y abordando las implicaciones socioeconómicas, de políticas viales y de privacidad \cite{kim2025unsupervised, wu2025visual}.

Consecuentemente, el objetivo principal es mapear y analizar estadísticamente las tendencias tecnológicas, entornos de validación y limitaciones de la visión artificial aplicada a ITS y monitoreo a partir de un corpus riguroso de 213 artículos científicos. Para guiar este análisis, planteamos responder tres preguntas de investigación: identificar las principales arquitecturas y enfoques metodológicos de IA (PI1); establecer las plataformas de hardware, software y entornos de validación predominantes (PI2); e identificar las limitaciones y retos críticos para delinear el trabajo futuro (PI3). El alcance se delimita a artículos de revistas indizadas y ponencias de congresos arbitrados publicados entre 2020 y 2026 en las bases de datos IEEE Xplore, ScienceDirect, SpringerLink y Scopus, excluyendo patentes, libros, reportes técnicos no arbitrados y trabajos puramente teóricos o sin validación empírica cuantitativa.

Finalmente, la estructura de este artículo se organiza de la siguiente manera. En la Sección II se detalla la metodología de revisión, describiendo el flujo PRISMA y el enfoque PICO. En la Sección III se presenta el análisis de resultados sobre los algoritmos, hardware y entornos experimentales. En la Sección IV se desarrolla la discusión de las implicaciones, la interpretabilidad y la adaptación de dominio. En la Sección V se examinan las limitaciones técnicas y los enfoques éticos de privacidad. Por último, en la Sección VI se exponen las conclusiones y directrices de trabajo futuro.
